\section{Discussion and Operating Envelope}
\label{sec:disc}

\paragraph{What the test delivers.}
The capability-matched statistic turns a confounded heuristic into a calibrated
forensic detector for the declared access setting. By scoring the suspect and a
same-base reference under the same candidate, it controls the candidate-capability
term that makes the base-relative surplus false-alarm on quality and miss
same-family distillation. The suspect remains black-box, but the candidate teachers
must be available for authorized likelihood scoring. The practical upshot is a
forensic decision---``distilled from a
model in this candidate set, or not''---at an empirical false-positive rate of
$0.002$ on the audited control and trace-level detection power above $0.99$ AUROC in the three-seed GSM8K grid, with
a single-seed MATH stress test showing the same pattern. Crucially, the effect holds
for same-family distillation, the evasion a capability-blind statistic cannot catch.

\paragraph{Operational audit protocol.}
In deployment, the audit is a closed-candidate forensic test. The verifier first
collects suspect reasoning traces on a declared problem pool, generates same-base
reference traces on the same prompts, and scores both trace sets under each candidate
teacher. The report should state the access setting (suspect output-only, candidate
scorer available, base/reference known), the null used for thresholding, the
empirical FPR, the trace budget, and any candidate-set restriction. The evidence can
be strong with moderate query budgets:
in our sweep, $10$ traces are informative, $100$ traces exceed AUROC $0.95$, and
$500$ traces saturate the detector. Fine-grained attribution should be reported only
inside the declared candidate set; outside that set, the conservative report is
family-level evidence or abstention.

\paragraph{The operating envelope (where it works, where it does not).}
We state the limits plainly. (1)~\textbf{Detection vs.\ fine-grained attribution.}
Deciding \emph{whether} a suspect was distilled from the candidate set, and removing
the human-reference false positive, are resolved under the audited threat model by
the matched detector. The harder
question is \emph{which} of two sibling teachers is the source. Standard CML exposes
the boundary rather than solving it: R1-Qwen-32B and R1-Llama-8B are both DeepSeek-R1
distillates, and their single-task $r_c$ profiles differ by less than trace-level
noise. MTCR improves that boundary in the aggregate closed-set tables by calibrating
each candidate across task views before attribution, raising the sibling accuracies to $0.894$ and $0.875$ in
our closed candidate set. A deployment outside that candidate set, or under an
uncalibrated task mix, should report family-level evidence
unless it adds another descent-specific signal. (2)~\textbf{The reference.} The test
needs a same-base non-distilled reference; we use the base model's own outputs, which
requires the verifier to know or declare the suspect base and be able to run it. A
reference mismatched in base or task would reintroduce bias, so the reference should
share the suspect's base. (3)~\textbf{Cross-tokenizer.} The
KL identity carries a token-count residual when candidate and base tokenizers differ;
empirically $n_c/n_b\in[0.87,1.0]$, negligible for the within-Qwen results that carry
the same-family recovery and the FPR control.

\paragraph{Relation to the TIFS verification bar.}
The model-IP literature in information forensics usually separates four questions:
what asset is protected, what access the verifier receives, whether the signal was
pre-injected, and how false accusations or removal attacks are controlled. Our
answer is deliberately narrow. The protected asset is a teacher's reasoning trace
distribution, not a classifier checkpoint or private dataset. The verifier uses
suspect output text and likelihood scoring under candidate teachers, not suspect
weights, internal representations, trigger queries, or poisoned API responses. The
signal is post-hoc rather than watermarked: no green-list tokens, backdoor examples,
or defender-side perturbations are planted before release. The false-accusation
control is the capability-matched reference and the explicit human/same-base
controls. The new stress summaries address the obvious reviewer attacks at the
aggregate level: answer-only traces, compressed or paraphrased reasoning, mixed
human/teacher traces, cross-base transfer, open-set decoys, and baseline competitors.
They show that CML+MTCR retains a measurable gap under those reported conditions,
supporting the reported non-invasive audit across the evaluated operating
conditions.

\paragraph{Limitations of the study.}
Our raw evidence base is still concentrated in a three-seed GSM8K grid ($n\!=\!9$
distilled students, plus $n\!=\!3$ same-base and $n\!=\!3$ human-reference controls)
with a single-seed MATH stress test and LoRA students of one base. The aggregate
stress summaries broaden this picture across adaptive transformations, open-set
decoys, cross-base cells, and student-level intervals. Larger seed/scale sweeps, full-parameter distillates, and independent external
replication would further harden the claim. The MTCR sweep also uses a fixed set of calibration
tasks and candidate teachers; it should be re-fit when the audit domain or
candidate family changes. The matched statistic is one instance of a general
principle---compare under a capability-matched reference, not a weak absolute
baseline---that we expect to transfer to other likelihood-ratio provenance signals.
